% TOP_PROBLEM: {"id": 30006768, "title": "Conformal Invariance of the Critical Random-Cluster Model with q <= 4", "category_id": 19, "category": {"id": 19, "name": "probability", "display_name": "Probability", "description": "Problems involving probability theory, stochastic processes, and random structures.", "slug": "probability", "order_index": 19, "created_at": "2026-07-31T15:26:25.671Z"}, "status": "open"}
% FIRST_POSTED: 2026-10-06
% ATTEMPT_STATUS: unresolved
% OPENAI_CLAIM: {"schema_version":1,"source_repo":"https://github.com/openai/math","source_commit":"adc7f1241b42e322a6451854ab7e4b4c146bf78a","release_date":"2026-10-06","match":"partial","resolution":"partial","independently_reviewed":false,"second_pass":"n/a","families":[{"family":"223","title":"Random-cluster interfaces: critical, disordered, thermal, and natural-time scaling","subject":"Probability and statistical mechanics","manuscripts":[{"title":"Conformal Limits of Critical Square-Lattice Random-Cluster Interfaces","pdf_path":"preprints/Conformal-Limits-of-Critical-Square-Lattice-Random-Cluster-Interfaces-October-5-2026/paper.pdf","readme_path":"preprints/Conformal-Limits-of-Critical-Square-Lattice-Random-Cluster-Interfaces-October-5-2026/README.md","date":"2026-10-05","theorem_ref":"intro:main (sections/introduction.tex:125)","statement_sources":[{"path":"preprints/Conformal-Limits-of-Critical-Square-Lattice-Random-Cluster-Interfaces-October-5-2026/build/sections/introduction.tex","line":125,"label":"intro:main"}]},{"title":"Natural Occupation Measures for Critical Square-Lattice FK Interfaces","pdf_path":"preprints/Natural-Occupation-Measures-for-Critical-Square-Lattice-FK-Interfaces-October-5-2026/natural-occupation-measures-critical-square-lattice-fk-interfaces.pdf","readme_path":"preprints/Natural-Occupation-Measures-for-Critical-Square-Lattice-FK-Interfaces-October-5-2026/README.md","date":"2026-10-05","theorem_ref":"thm:main (sections/01-introduction.tex:114)","statement_sources":[{"path":"preprints/Natural-Occupation-Measures-for-Critical-Square-Lattice-FK-Interfaces-October-5-2026/build/sections/01-introduction.tex","line":114,"label":"thm:main"}]},{"title":"Self-dual random-cluster interfaces below one","pdf_path":"preprints/Self-dual-random-cluster-interfaces-below-one-September-23-2026/paper.pdf","readme_path":"preprints/Self-dual-random-cluster-interfaces-below-one-September-23-2026/README.md","date":"2026-09-23","theorem_ref":"thm:main (sections/00-introduction.tex:83)","statement_sources":[{"path":"preprints/Self-dual-random-cluster-interfaces-below-one-September-23-2026/build/sections/00-introduction.tex","line":83,"label":"thm:main"}]},{"title":"Square-lattice FK interfaces and nested loops for 1≤q<4","pdf_path":"preprints/Square-lattice-FK-interfaces-and-nested-loops-for-1-leq-q-lt-4-September-23-2026/paper.pdf","readme_path":"preprints/Square-lattice-FK-interfaces-and-nested-loops-for-1-leq-q-lt-4-September-23-2026/README.md","date":"2026-09-23","theorem_ref":"intro:main (sections/introduction.tex:193)","statement_sources":[{"path":"preprints/Square-lattice-FK-interfaces-and-nested-loops-for-1-leq-q-lt-4-September-23-2026/build/sections/introduction.tex","line":193,"label":"intro:main"}]}],"lean":null,"reasoning_trace":null}]}
\documentclass[11pt]{article}
\usepackage[margin=25mm]{geometry}
\usepackage{fontspec}
\setmainfont{Latin Modern Roman}
\usepackage{amsmath,amssymb,mathtools}
\usepackage{unicode-math}
\setmathfont{Latin Modern Math}
\usepackage{enumitem}
\usepackage{xurl,hyperref,bookmark}
\hypersetup{colorlinks=true,urlcolor=blue}
\setcounter{secnumdepth}{0}
\setlength{\parindent}{0pt}
\setlength{\parskip}{5pt}
\setlength{\emergencystretch}{3em}
\begin{document}
\section*{Conformal Invariance of the Critical Random-Cluster Model with q <= 4}

\subsection{OpenAI's claim}
The linked manuscripts present the following claims by OpenAI: The later interface theorem extends the square-lattice FK claim through $q=4$, with the same SLE parameter relation and a separate complete nested-loop limit; For $1\leq q<4$, the rescaled edge-counting measure of the critical square FK interface is claimed to converge with the curve to its SLE Minkowski-content measure; For $0<q<1$, the self-dual square-lattice FK interface is claimed to converge to chordal SLE with $\kappa=4\pi/\arccos(-\sqrt q/2)$ in smooth Jordan domains; For $1\leq q<4$, critical square-lattice FK interfaces are claimed to converge as curves to the predicted SLE, with a separate convergence theorem for the complete nested plane CLE collection. Together they provide the following support within the scopes described below.

\subsection{Mathematical statement of the claimed result}
\paragraph{Result 1.} Fix \(q\in[1,4]\), \(p_q=\sqrt q/(1+\sqrt q)\), and
\[
 \kappa(q)=\frac{4\pi}{\arccos(-\sqrt q/2)}.
\]
Use the square-lattice random-cluster law with edge-subset weights \(p_q^{|A|}(1-p_q)^{|E|-|A|}q^{k_\xi(A)}\), where \(\xi\) wires one boundary arc and leaves the other free. For a bounded Jordan domain with two marked boundary points, approximate its boundary uniformly by simple lattice polygons, preserving the marks. The claim gives
\[
 \eta_n\ \Longrightarrow\ \operatorname{SLE}_{\kappa(q)}
\]
for the Dobrushin interface in uniform curve distance modulo increasing reparametrization. Under the free infinite-volume law, its full nested loop collection satisfies
\[
 \Gamma_\delta\ \Longrightarrow\ \text{nested whole-plane }\operatorname{CLE}_{\kappa(q)}
\]
in the spherical matching topology: all loops above a shrinking diameter threshold are matched with shrinking uniform spherical error. Both limits hold along every mesh sequence, with \(q\) fixed.

\paragraph{Result 2.} Fix \(q\in[1,4)\), \(p_q=\sqrt q/(1+\sqrt q)\), \(\kappa=4\pi/\arccos(-\sqrt q/2)\), and \(d=1+\kappa/8\). On mesh \(1/n\) in \(D=(0,1)^2\), use critical FK weights and wired/free boundary arcs between \((0,0)\) and \((1,1)\). Let \(\eta_n\) be the oriented Dobrushin strand, and \(z_{n,j}\) its medial-edge traversal midpoints, counting boundary traversals but not terminal half-edges. With \(\Pi\) the coordinatewise projection onto \([0,1]^2\), the source claims a constant \(c(q)>0\) such that
\[
 ([\eta_n],c(q)n^{-d}\sum_{j=1}^{N_n}\delta_{\Pi(z_{n,j})})
 \Longrightarrow([\eta],\mu_\eta).
\]
Here \(\eta\) is chordal \(\operatorname{SLE}_\kappa\), curves retain order modulo increasing reparametrization, and measures converge weakly on \(\overline D\). The limiting content measure is the weak limit of \(r^{d-2}\Pi_*[\mathbf1_{\{\operatorname{dist}(z,\eta)<r\}}dA(z)]\), using whole-plane area. The two limit coordinates belong to the same random curve.

\paragraph{Result 3.} For fixed \(q\in(0,1)\), use square-lattice FK edge-subset weights at \(p=\sqrt q/(1+\sqrt q)\), with one boundary arc wired and the other free. Choose a bounded domain \(D\) whose Jordan boundary is smooth, and choose different marked boundary points \(a,b\). Approximate it by simple closed nearest-neighbor polygons whose boundary parametrizations converge uniformly, and whose marked cyclic parameters converge. The source claims
\[
 \eta_\delta\Longrightarrow\operatorname{SLE}_\kappa(D;a,b),\qquad
 \kappa=\frac{4\pi}{\arccos(-\sqrt q/2)}\in(6,8),
\]
for the Dobrushin medial strand, along the entire mesh sequence, in uniform curve distance modulo increasing reparametrization. The theorem's domain assumption is smooth boundary, rather than every Jordan domain.

\paragraph{Result 4.} Fix \(1\le q<4\), \(p_q=\sqrt q/(1+\sqrt q)\), and \(\kappa(q)=4\pi/\arccos(-\sqrt q/2)\). The square-lattice critical random-cluster model is claimed to satisfy two limits: its wired/free Dobrushin strand in any marked bounded Jordan domain converges to chordal \(\operatorname{SLE}_{\kappa(q)}\); the complete interface loop collection under the free infinite-volume law converges to nested whole-plane \(\operatorname{CLE}_{\kappa(q)}\). The first uses uniformly convergent marked simple-polygon boundary parametrizations and uniform curve distance modulo increasing reparametrization. The second uses spherical matching of all loops of diameter exceeding a shrinking threshold, retaining every nesting level. The result is for fixed \(q<4\); it asserts no uniformity as \(q\to4\).

\subsection{Outline of the argument}
The tiled comparison and angular-transfer method is supplemented by endpoint estimates at $q=4$. Boundary current tests identify the driving process, while full exploration and localization retain macroscopic loops; at $q=1$ this also yields the bond-percolation crossing formula. Row slicing and angular electric spectra determine Coulomb correlations with quantitative remainders. A two-arm outward coupling fixes the single-edge amplitude, and integrated correlation bounds identify the occupation measure and its total mass with one deterministic normalization.

Finite connection and boundary-partition comparisons replace positive-association arguments that fail below one. Sector decay localizes irregular tiled disks, and angular transfer gives a boundary observable determining the Loewner law along the full mesh sequence. Positive tiled laws and short-distance comparisons localize the interface in irregular disks. Angular transfer modes and a four-boundary-change current observable identify the Loewner driver; exploration iteration then matches the nested loops, retaining their multiplicities and traversal structure.

\subsection{Scope relative to this problem}
The interface theorem proves full-sequence Dobrushin SLE and nested CLE convergence for 1<=q<=4 on square-lattice Jordan approximations; the problem statement target also includes 0<q<1, so this paper alone does not cover the full parameter range. The cited theorem proves joint curve/occupation-measure convergence for critical square-lattice FK interfaces with 1<=q<4; forgetting the measure gives that subrange of our interface target, but not q<1 or q=4. The cited theorem proves the square-lattice Dobrushin FK interface SLE limit for each 0<q<1; our range includes 1<=q<=4, which this individual manuscript does not prove. The interface theorem proves chordal SLE and nested CLE limits for the critical square-lattice FK model only for 1<=q<4; our intended FK range is 0<q<=4, so q<1 and the endpoint q=4 are omitted at manuscript level. This entry records partial support; complete coverage of the question is not established by this catalogue.

\subsection{Formal verification}
Family 223 has no Lean scope page in this release. Its selected result is not presented as formalized. This site has not executed the Lean code. The scope descriptions do not certify the remaining manuscript arguments.

\subsection{Status caveat}
OpenAI attributes these manuscripts to work by an internal model. Its release cautions that arguments without formalization can contain mistakes. This catalogue has not independently checked the proofs; the labels report the mathematical claims attributed to OpenAI.

\subsection{Sources}
\begin{itemize}
\item OpenAI, \emph{Conformal Limits of Critical Square-Lattice Random-Cluster Interfaces}, 2026-10-05; citation key \texttt{OAI:Conformal-Limits-of-Critical-Square-Lattice-Random-Cluster-Interfaces-October-5-2026}. \href{https://github.com/openai/math/blob/adc7f1241b42e322a6451854ab7e4b4c146bf78a/preprints/Conformal-Limits-of-Critical-Square-Lattice-Random-Cluster-Interfaces-October-5-2026/paper.pdf}{PDF (GitHub)}; \href{https://github.com/openai/math/raw/adc7f1241b42e322a6451854ab7e4b4c146bf78a/preprints/Conformal-Limits-of-Critical-Square-Lattice-Random-Cluster-Interfaces-October-5-2026/paper.pdf}{PDF (direct)}; \href{https://github.com/openai/math/blob/adc7f1241b42e322a6451854ab7e4b4c146bf78a/preprints/Conformal-Limits-of-Critical-Square-Lattice-Random-Cluster-Interfaces-October-5-2026/README.md}{Manuscript page and citation}; \href{https://github.com/openai/math/blob/main/preprints/Conformal-Limits-of-Critical-Square-Lattice-Random-Cluster-Interfaces-October-5-2026/paper.pdf}{Latest version}; Statement source: \href{https://github.com/openai/math/blob/adc7f1241b42e322a6451854ab7e4b4c146bf78a/preprints/Conformal-Limits-of-Critical-Square-Lattice-Random-Cluster-Interfaces-October-5-2026/build/sections/introduction.tex\#L125}{intro:main}.
\item OpenAI, \emph{Natural Occupation Measures for Critical Square-Lattice FK Interfaces}, 2026-10-05; citation key \texttt{OAI:Natural-Occupation-Measures-for-Critical-Square-Lattice-FK-Interfaces-October-5-2026}. \href{https://github.com/openai/math/blob/adc7f1241b42e322a6451854ab7e4b4c146bf78a/preprints/Natural-Occupation-Measures-for-Critical-Square-Lattice-FK-Interfaces-October-5-2026/natural-occupation-measures-critical-square-lattice-fk-interfaces.pdf}{PDF (GitHub)}; \href{https://github.com/openai/math/raw/adc7f1241b42e322a6451854ab7e4b4c146bf78a/preprints/Natural-Occupation-Measures-for-Critical-Square-Lattice-FK-Interfaces-October-5-2026/natural-occupation-measures-critical-square-lattice-fk-interfaces.pdf}{PDF (direct)}; \href{https://github.com/openai/math/blob/adc7f1241b42e322a6451854ab7e4b4c146bf78a/preprints/Natural-Occupation-Measures-for-Critical-Square-Lattice-FK-Interfaces-October-5-2026/README.md}{Manuscript page and citation}; \href{https://github.com/openai/math/blob/main/preprints/Natural-Occupation-Measures-for-Critical-Square-Lattice-FK-Interfaces-October-5-2026/natural-occupation-measures-critical-square-lattice-fk-interfaces.pdf}{Latest version}; Statement source: \href{https://github.com/openai/math/blob/adc7f1241b42e322a6451854ab7e4b4c146bf78a/preprints/Natural-Occupation-Measures-for-Critical-Square-Lattice-FK-Interfaces-October-5-2026/build/sections/01-introduction.tex\#L114}{thm:main}.
\item OpenAI, \emph{Self-dual random-cluster interfaces below one}, 2026-09-23; citation key \texttt{OAI:Self-dual-random-cluster-interfaces-below-one-September-23-2026}. \href{https://github.com/openai/math/blob/adc7f1241b42e322a6451854ab7e4b4c146bf78a/preprints/Self-dual-random-cluster-interfaces-below-one-September-23-2026/paper.pdf}{PDF (GitHub)}; \href{https://github.com/openai/math/raw/adc7f1241b42e322a6451854ab7e4b4c146bf78a/preprints/Self-dual-random-cluster-interfaces-below-one-September-23-2026/paper.pdf}{PDF (direct)}; \href{https://github.com/openai/math/blob/adc7f1241b42e322a6451854ab7e4b4c146bf78a/preprints/Self-dual-random-cluster-interfaces-below-one-September-23-2026/README.md}{Manuscript page and citation}; \href{https://github.com/openai/math/blob/main/preprints/Self-dual-random-cluster-interfaces-below-one-September-23-2026/paper.pdf}{Latest version}; Statement source: \href{https://github.com/openai/math/blob/adc7f1241b42e322a6451854ab7e4b4c146bf78a/preprints/Self-dual-random-cluster-interfaces-below-one-September-23-2026/build/sections/00-introduction.tex\#L83}{thm:main}.
\item OpenAI, \emph{Square-lattice FK interfaces and nested loops for 1\(\leq\)q<4}, 2026-09-23; citation key \texttt{OAI:Square-lattice-FK-interfaces-and-nested-loops-for-1-leq-q-lt-4-September-23-2026}. \href{https://github.com/openai/math/blob/adc7f1241b42e322a6451854ab7e4b4c146bf78a/preprints/Square-lattice-FK-interfaces-and-nested-loops-for-1-leq-q-lt-4-September-23-2026/paper.pdf}{PDF (GitHub)}; \href{https://github.com/openai/math/raw/adc7f1241b42e322a6451854ab7e4b4c146bf78a/preprints/Square-lattice-FK-interfaces-and-nested-loops-for-1-leq-q-lt-4-September-23-2026/paper.pdf}{PDF (direct)}; \href{https://github.com/openai/math/blob/adc7f1241b42e322a6451854ab7e4b4c146bf78a/preprints/Square-lattice-FK-interfaces-and-nested-loops-for-1-leq-q-lt-4-September-23-2026/README.md}{Manuscript page and citation}; \href{https://github.com/openai/math/blob/main/preprints/Square-lattice-FK-interfaces-and-nested-loops-for-1-leq-q-lt-4-September-23-2026/paper.pdf}{Latest version}; Statement source: \href{https://github.com/openai/math/blob/adc7f1241b42e322a6451854ab7e4b4c146bf78a/preprints/Square-lattice-FK-interfaces-and-nested-loops-for-1-leq-q-lt-4-September-23-2026/build/sections/introduction.tex\#L193}{intro:main}.
\item \href{https://github.com/openai/math/blob/adc7f1241b42e322a6451854ab7e4b4c146bf78a/CONTENTS.md}{OpenAI catalogue entry}.
\end{itemize}
\end{document}
